\section{Thiết kế Edge Agent}

Edge Agent là tiến trình thường trực duy nhất được cài tĩnh trên Edge Node; mọi
AI Runtime do Agent tạo động từ image runtime. Các thành phần chính được thể hiện
trong Hình~\ref{fig:edge-agent-components}.

\begin{figure}[H]
\centering
\resizebox{0.95\textwidth}{!}{%
\begin{tikzpicture}[font=\small,
  b/.style={draw, rounded corners, align=center, minimum height=1.0cm, text width=3.9cm, fill=blue!5}]
  \node[b, fill=purple!6] (mqtt) at (0,4.4) {MQTT\\(command/+)};
  \node[b, fill=green!6] (api) at (4.4,4.4) {Local API :8081\\(REST)};
  \node[b] (cmd) at (2.2,2.6) {CommandHandler\\(chống lệnh trùng)};
  \node[b] (rt) at (2.2,0.6) {RuntimeService\\(apply / reconcile)};
  \node[b] (desired) at (-3.2,0.6) {JsonDesiredState\\Repository};
  \node[b] (drv) at (-0.2,-1.6) {DockerRuntimeDriver\\(Docker CLI)};
  \node[b] (fetch) at (4.4,-1.6) {HttpArtifactFetcher\\(cache model)};
  \node[b] (agent) at (8.2,0.6) {AgentService\\(vòng lặp định kỳ)};
  \node[b] (state) at (10.0,2.6) {StateReportService\\(10 s)};
  \node[b] (tele) at (10.0,-1.6) {TelemetryReportService\\(5 s)};
  \draw[flowarrow] (mqtt) -- (cmd);
  \draw[flowarrow] (api) -- (cmd);
  \draw[flowarrow] (cmd) -- (rt);
  \draw[flowarrow] (rt) -- (drv);
  \draw[flowarrow] (rt) -- (fetch);
  \draw[flowarrow, <->] (rt) -- (desired);
  \draw[flowarrow] (agent) -- node[font=\scriptsize, above=3pt]{reconcile 15 s} (rt);
  \draw[flowarrow] (agent) -- (state);
  \draw[flowarrow] (agent) -- (tele);
\end{tikzpicture}}
\caption{Các thành phần chính của Edge Agent}
\label{fig:edge-agent-components}
\end{figure}

\subsection{Vòng lặp điều khiển}

Khi khởi động, AgentService kết nối broker, khôi phục desired state từ đĩa,
đăng ký topic lệnh và phát ngay một NodeState. Sau đó Agent chạy ba vòng lặp độc
lập (chu kỳ cấu hình được):

\begin{itemize}
  \item \textbf{State} (10~s): phát NodeState dạng retained, đóng vai trò heartbeat
        và báo cáo toàn bộ runtime. Agent ở trạng thái DEGRADED nếu có runtime
        mong muốn chạy nhưng đang FAILED hoặc DEGRADED.
  \item \textbf{Telemetry} (5~s): thu và phát telemetry \texttt{system} và \texttt{gpu}.
  \item \textbf{Reconcile} (15~s): so sánh desired state với trạng thái container
        thực tế và hội tụ.
\end{itemize}

Khi dừng, Agent hủy các vòng lặp và phát NodeState OFFLINE trước khi ngắt kết nối.

\subsection{Xử lý lệnh}

CommandHandler nhận lệnh từ MQTT hoặc từ API cục bộ; hai nguồn đi qua cùng một
đường xử lý nên kết quả đều được báo về Cloud. Lệnh không đúng schema bị từ chối;
lệnh có \texttt{command\_id} đã xử lý (bộ nhớ đệm 256 id gần nhất) bị bỏ qua để
chống giao lặp của broker. Sau khi thực thi, Agent phát \texttt{DeploymentEvent}
chứa \texttt{command\_id}, \texttt{deployment\_id}, trạng thái deployment và trạng
thái runtime, rồi phát ngay một NodeState mới. Lỗi nghiệp vụ (\texttt{AppException})
và lỗi bất ngờ đều được chuyển thành kết quả FAILED có thông điệp, không làm dừng Agent.

\subsection{Desired state và reconcile}

RuntimeService giữ desired state của mọi runtime trên node và là nơi duy nhất
thay đổi container. Mọi thao tác đi qua một khóa bất đồng bộ nên các lệnh trên
cùng node được tuần tự hóa. Các thao tác đều idempotent:

\begin{itemize}
  \item \texttt{apply(spec)}: nếu spec giống hệt và runtime đang RUNNING thì không
        làm gì; nếu khác, hủy container cũ, lưu desired state và khởi chạy mới.
  \item \texttt{start}/\texttt{stop}: đặt \texttt{desired\_running} rồi thao tác
        container; \texttt{remove} với runtime không tồn tại vẫn thành công.
  \item \texttt{reconcile()}: với runtime mong muốn dừng mà container đang chạy thì
        dừng; mong muốn chạy mà container mất hoặc đã thoát thì ghi lý do (mã thoát,
        lỗi), xóa container cũ và tạo lại; sau \texttt{max\_failures} (5) lần thất bại
        liên tiếp thì đánh dấu FAILED và ngừng thử.
  \item \texttt{restore()}: sau khi Agent khởi động lại, đọc desired state từ đĩa và
        tìm container còn sống theo nhãn \texttt{edgeops.runtime\_id} để tiếp quản
        thay vì tạo trùng.
\end{itemize}

Mỗi lần trạng thái runtime thay đổi, Agent phát một \texttt{EdgeEvent}
(\texttt{runtime.status}); chuyển sang FAILED được phát trên kênh \texttt{error}
với mức ERROR.

\subsection{Khởi chạy runtime}

Quy trình \texttt{\_launch} gồm: (1) chuyển CREATING; (2) nếu spec có
\texttt{model\_uri}, tải model qua ArtifactFetcher và gắn vào container tại
\path{/models/<file>} ở chế độ chỉ đọc; (3) pull image theo chính sách
\texttt{always}/\texttt{missing}/\texttt{never}; (4) chuyển STARTING và chạy
container với biến môi trường \texttt{RUNTIME\_ID}, \texttt{RUNTIME\_NAME},
\texttt{RUNTIME\_CONFIG}, \texttt{MODEL\_PATH}, \texttt{NODE\_ID},
\texttt{AGENT\_URL}; (5) chuyển RUNNING. Giới hạn tài nguyên trong spec được ánh
xạ sang tham số Docker: \texttt{--gpus}, \texttt{--runtime nvidia} (Jetson),
\texttt{--cpus}, \texttt{--memory}, \texttt{--device} (camera). Container được gắn
chính sách \texttt{--restart unless-stopped} nên Docker tự khởi động lại khi
thiết bị reboot.

\subsection{API cục bộ}

Agent cung cấp REST API trên cổng 8081 (bảo vệ bằng API key) gồm hai nhóm:
nhóm quản lý (\texttt{GET /status}, \texttt{GET/POST/PUT/DELETE /runtimes},
\texttt{POST /runtimes/\{id\}/\{start|stop|restart\}}) cho phép vận hành tại chỗ
khi không có Cloud; và nhóm báo cáo (\texttt{POST /runtimes/\{id\}/telemetry},
\texttt{POST /runtimes/\{id\}/events}) để AI Runtime gửi chỉ số. Agent làm giàu mẫu
telemetry runtime với \texttt{node\_id}, trạng thái, số lần restart, tính
\texttt{effective\_fps\_ratio} = FPS thực/FPS mục tiêu, rồi chuyển tiếp lên Cloud.
